% SOURCE_STATEMENT: "For rate-independent fracture and damage models coupled with rate-dependent processes, are there efficient feasible numerical schemes that reliably approximate the intended solution class, and what mathematical and mechanical framework should describe the coupling?"
% Preparation and final self-audit: CONTRIBUTING.md prompts C and D; source/audit artifacts retained in the local physics_statements_2026_10_11 batch.
% !TeX program = lualatex
\documentclass[11pt]{article}
\usepackage[margin=25mm]{geometry}
\usepackage{fontspec}
\setmainfont{Latin Modern Roman}
\usepackage{amsmath,amssymb,mathtools}
\usepackage{unicode-math}
\setmathfont{Latin Modern Math}
\usepackage{xurl,hyperref}
\hypersetup{colorlinks=true,urlcolor=blue}
\setcounter{secnumdepth}{0}
\setlength{\parindent}{0pt}
\setlength{\parskip}{5pt}
\setlength{\emergencystretch}{3em}
\newcommand{\sref}[2]{\hyperlink{source:#1:#2}{[S#2]}}
\newcommand{\cataloguescope}[1]{\paragraph{Recorded target.}#1}
\newcommand{\R}{\mathbb R}
\newcommand{\C}{\mathbb C}
\newcommand{\N}{\mathbb N}
\newcommand{\Z}{\mathbb Z}
\newcommand{\Q}{\mathbb Q}
\newcommand{\F}{\mathbb F}
\newcommand{\eps}{\varepsilon}
\begin{document}
\section{Numerical Schemes for Coupled Fracture and Dissipation}\label{statement}

\subsection{Definitions and mathematical statement}

\paragraph{Shared notation.}$\mathbb N=\{1,2,\ldots\}$; the standard number fields have their usual meanings. All additional conventions are stated locally.


\paragraph{Definitions and conventions.}
A rate-independent model has dissipation positively homogeneous of degree one in the evolution rate, so its response is invariant under suitable time reparametrization. Rate-dependent processes have a distinct time scale, often through viscous dissipation. Fracture and damage variables obey irreversibility constraints. Energetic, local and vanishing-viscosity solution concepts may differ at jumps when the energy is nonconvex. A reliable numerical scheme must identify and converge to the intended concept while resolving the coupling.

\paragraph{Target retained from the canonical source catalogue.}
\paragraph{Catalogue wording retained for traceability.}
For rate-independent fracture and damage models coupled with rate-dependent processes, are there efficient feasible numerical schemes that reliably approximate the intended solution class, and what mathematical and mechanical framework should describe the coupling?

\paragraph{Source scope and qualification.}
The desired solution class must be selected explicitly: the report emphasizes that energetic and vanishing-viscosity evolutions can have different jumps. Efficiency alone is not convergence to the intended physical and mathematical evolution.
\subsection{Short English statement}
Develop a sound coupled framework and practical convergent schemes for fracture or damage interacting with time-dependent dissipation.
\subsection{Sources}
\begin{itemize}
\item[\hypertarget{source:statement:1}{S1}] ULAM.AI and the UnsolvedMath Contributors, UnsolvedMath, catalogue entry. Catalogue curation under CC BY 4.0; definitions and English exposition added locally.
\url{https://huggingface.co/datasets/ulamai/UnsolvedMath}.
\item[\hypertarget{source:statement:2}{S2}] Mechanics of Materials: Towards Predictive Methods for Kinetics in Plasticity, Fracture, and Damage (2020), Oberwolfach Reports, DOI 10.4171/owr/2020/13, pp. 713--783. Primary report accessed 2026-10-11; relevant PDF pages 35.
\url{https://doi.org/10.4171/owr/2020/13}.
\item[\hypertarget{source:statement:3}{S3}] Rossi--Thomas (2017) (catalogue-listed context; not independently consulted in this preparation).
\url{https://doi.org/10.1137/15M1051567}.
\end{itemize}
\label{end:statement}
\end{document}
